\begin{comment}
% \vspace{-0.2in}
\section{Related Work}\label{sec:related}
% \vspace{-0.1in}
% \paragraph{Implicit Bias of SGD.}  
% Several studies have characterized the bias induced by multi‐pass SGD.  \citet{haochen2021shape} showed that SGD with label noise recovers sparse ground truth in quadratic models.  \citet{blanc2020implicit} proved that, under MSE and label noise, SGD’s fixed point minimizes the Hessian trace on the zero‐loss manifold; \citet{damian2021label} extended this to global convergence and large learning rates.  \citet{li2021happens} introduced an SDE framework for SGD’s bias near the minimizer manifold, later refined by \citet{gu2023andwhendoeslocal} to Local SGD~\citep{lin2018don}, and by \citet{wang2023marginal} to include momentum’s negligible effect.

% Parallel work on \emph{implicit gradient regularization} (IGR) derives higher‐order terms for full‐batch GD \citep{barrett2020implicit} and extends to Adam \citep{cattaneo2024implicitbiasadam,cattaneo2025memoryoptimizationalgorithmsimplicitly}.  While \citet{cattaneo2024implicitbiasadam} argued that full-batch Adam with constant learning rate approximately follows an ODE, whcih implies an anti‐regularizes sharpness when $\beta_1<\beta_2$, our $\O(\eta^{-2})$‐time SDE analysis shows the empirical Adam can be characterized by an slow SDE in a long horizon, which shows Adam still regularizes sharpness in a long horizon, offering a complementary perspective.
% \xinghan{!!!vfk look at this $\sim$}
\paragraph{Implicit Bias of Adam.}  
Despite Adam’s widespread use, its implicit bias remains underexplored.  \citet{qian2019implicit} and \citet{xie2024implicit} analyzed AdaGrad and AdamW, but these techniques do not apply directly to Adam.  \citet{wang2021implicit} showed Adam’s regularizer matches SGD’s under restrictive gradient‐magnitude assumptions, and \citet{zhang2024implicitbiasadamseparable} treated only linearly separable data, limiting practical relevance.

Also, another line of works on \emph{implicit gradient regularization} (IGR) derive higher‐order approximations for full‐batch GD \citep{barrett2020implicit} and extend to Adam \citep{cattaneo2024implicitbiasadam,cattaneo2025memoryoptimizationalgorithmsimplicitly}.  In particular, \citet{cattaneo2024implicitbiasadam} argued that full-batch Adam with constant learning rate approximately follows an ODE that anti‐regularizes sharpness when $\beta_1<\beta_2$. Our work analyzes the dynamics of Adam for $\O(\eta^{-2})$ steps, a longer horizon than \citet{cattaneo2024implicitbiasadam}. Our analysis shows that with gradient noise, Adam can be characterized by an slow SDE that regularizes sharpness in the long term, offering a complementary perspective.

\paragraph{Slow SDE Approximation.}  
To capture long‐term behavior, we adopt the \emph{slow SDE} technique of \citet{li2021happens} and \citet{gu2023andwhendoeslocal}.  Standard SDE approximations (\citealp{li2018stochasticmodifiedequationsdynamics,li2021validitymodelingsgdstochastic,cattaneo2024implicitbiasadam,malladi2024sdesscalingrulesadaptive}) focus on the $\tilde{\O}(\eta^{-1})$ convergence phase and fail on the manifold.  In contrast, slow SDEs peel off convergence to track the $\O(\eta^{-2})$ manifold dynamics accurately.

For more detailed discussion on the related work, please refer to \Aref{sec:add-related}.
% \paragraph{Adaptive Gradient Methods.}  
% Since Adam’s introduction~\citep{kingma2014adam}, adaptive methods (AdamW \citep{loshchilov2017decoupled}, AdaFactor \citep{shazeer2018adafactor}, Adam-mini \citep{zhang2025adamminiusefewerlearning}, Adalayer \citep{zhao2025deconstructingmakesgoodoptimizer}, AdaSGD \citep{wang2020adasgd}) and earlier schemes (RMSprop \citep{hinton2012neural}) have dominated practice, offering stability across domains.  However, their implicit biases and long‐term dynamics remain theoretically unresolved.
\end{comment}

\section{Related Work}\label{sec:related}

\paragraph{Implicit Bias of SGD.}
% The same tool has also been applied to LocalSGD \citep{lin2018don}, a distributed variant of SGD, by \citet{gu2023andwhendoeslocal}.
%The implicit bias of SGD has been studied over time; 


A line of works studies the implicit bias of SGD when training overparameterized models with label noise.
\citet{blanc2020implicit} proved that with $\ell_2$ loss and label noise, with high probability, SGD will move away from points on the minimizer manifold $\Gamma$ if and only if $\tr(\mH)$ is not locally minimal within $\Gamma$. \citet{haochen2021shape} further extended this local characterization to a global convergence result under the diagonal net setting \citep{woodworth2020kernel}.
% , and also emphasized that this bias relies specifically on label noise, whereas SGD with Gaussian noise or plain Gradient Descent (GD) fail to induce the same regularization effect. 
However, \citet{haochen2021shape} relied on a manually designed, non-constant learning-rate schedule. \citet{damian2021label} overcame this issue by proving that the same implicit bias holds under constant learning rates and for any smooth loss satisfying the Kurdyka–Łojasiewicz condition. They also proved the same implicit bias for SGD with momentum (SGDM).
% and their framework accommodates general noise covariance matrices. 

Up to this point, all analyses were not able to track the optimization over $\O(\eta^{-2})$ iterations, but this is necessary to capture the entire sharpness reduction process after converging to the minimizer manifold.
\citet{li2021happens} were the first to tackle this problem, which provided an accurate SDE-based characterization of the long-term behavior of SGD over $\O(\eta^{-2})$ time. 
Their SDE characterization also goes beyond the label noise case and is able to capture the implicit bias induced by general forms of gradient noise.
However, the SDE derivation in \citet{li2021happens} is specific to plain SGD and cannot be directly extended to other optimizers. \citet{gu2023andwhendoeslocal} later termed this form of SDE \emph{the slow SDE} and derived it for local SGD~\citep{lin2018don} in a way that is potentially generalizable to a broader class of optimizers. 
\citet{wang2023marginal} reinforced that momentum does not alter the implicit bias by showing that SGD and SGDM are characterized by the same slow SDE.

% \xinghan{IGR under construction}

% \citet{li2018stochasticmodifiedequationsdynamics}, who derived from the stochastic modified equation of SGD that its dynamics correspond to a perturbed gradient flow with an additional gradient-norm regularizer. This concept was further formalized and 

Another line of works characterizes the implicit bias of SGD through the lens of \textit{Gradient Regularization (GR)}. \citet{li2018stochasticmodifiedequationsdynamics} first derived a \emph{stochastic modified equation} of SGD, showing that introducing an additional gradient norm regularizer to the drift term of SDE helps approximate SGD more accurately. Later \citet{barrett2020implicit} also discovered the same regularizer that can be added to gradient flow to approximate GD in a higher order, and termed this behavior the \textit{Implicit Gradient Regularization (IGR)}.
\citet{smith2021originimplicitregularizationstochastic} specifically studied and highlighted the IGR for SGD. 
Empirical studies \citep{geiping2022stochastictrainingnecessarygeneralization, novack2022disentanglingmechanismsimplicitregularization} showed that GR is not necessarily implicit by demonstrating that explicitly regularizing gradient norm for SGD with larger or even full batch size can recover the generalization performance of small-batch SGD. GR was also found by \citet{karakida2023understandinggradientregularizationdeep} to be connected with Sharpness-Aware Minimization (SAM) \citep{foret2021sharpnessaware} and inspired improvements upon SAM in terms of generalization performance \citep{zhao2022penalizinggradientnormefficiently}. 
The IGR method was later generalized to the analysis of as Adam and AdamW by \citet{cattaneo2024implicitbiasadam} and \citet{ cattaneo2025memoryoptimizationalgorithmsimplicitly}. We further discuss our relationship to \citet{cattaneo2024implicitbiasadam} in the next paragraph.


% Another line of works characterizes the implicit bias of SGD through the lens of \textit{Implicit Gradient Regularization (IGR)}. For full-batch GD, this method derives higher-order correction terms to the gradient-flow ODE to better approximate discrete GD iterates \citep{barrett2020implicit}; This idea was later generalized to the analysis of Adam by \citet{cattaneo2024implicitbiasadam} and further to AdamW by \citet{cattaneo2025memoryoptimizationalgorithmsimplicitly}. We further discuss our relationship to \citet{cattaneo2024implicitbiasadam} in the next paragraph.

\paragraph{Implicit Bias of Adam.}
On the theoretical side, the current literature still lacks a rigorous understanding of the implicit bias of Adam, although it has been used more widely than SGD in practical deep learning, especially for large language models. However, many efforts have been made on this 
problem. \citet{qian2019implicit} and \citet{xie2024implicit} characterized the implicit biases of AdaGrad and AdamW, respectively. However, their methods do not extend to Adam. \citet{wang2021implicit} showed that Adam's implicit regularizer is identical to that of SGD, while their result requires that the gradient coordinates have magnitude at most $\epsilon$, which is typically not feasible in practice. \citet{zhang2024implicitbiasadamseparable}'s analysis is also limited in scope, as it focuses on Adam's implicit bias on linearly separable data, a condition generally not met by real-world applications. A work using IGR to analyze the implicit bias of Adam, \citet{cattaneo2024implicitbiasadam}, seems similar at first glance to our study, but we are actually offering a different angle on the implicit bias of Adam. In particular, \citet{cattaneo2024implicitbiasadam} argued that full-batch Adam with constant learning rate approximately follows an ODE that anti‐regularizes sharpness when $\beta_1<\beta_2$. Our work analyzes the dynamics of Adam for $\O(\eta^{-2})$ steps, a longer horizon than \citet{cattaneo2024implicitbiasadam}. Our analysis shows that with gradient noise, Adam can be characterized by a slow SDE that regularizes sharpness in the long term, offering a complementary perspective.


\paragraph{Approximation of Stochastic First-Order Optimizers with SDE.} 
Optimizers such as SGD and Adam take $\tilde{\O}(\eta^{-1})$ steps to converge onto the manifold, and then move along the manifold for $\O(\eta^{-2})$ steps, during which the optimizer will be dominated by a slower implicit regularization dynamics, different from the convergence dynamics earlier. Before slow SDE was introduced by \citet{li2021happens}, the conventional SDE approximations \citep{li2018stochasticmodifiedequationsdynamics,li2021validitymodelingsgdstochastic,cattaneo2024implicitbiasadam,malladi2024sdesscalingrulesadaptive} track the iteration during the convergence phase, but they struggle to bound the approximation error if extended to the manifold phase. Slow SDE tackles this problem by peeling the convergence dynamics off and approximating the iteration's projection on the manifold only. In this way, slow SDE manages to track iteration throughout the manifold phase for $\O(\eta^{-2})$ time. This idea will be made explicit in \Cref{sec:prelim}. 
In this work, to fill in the gaps and provide a theoretical analysis that tracks iterations of Adam for a sufficiently long time, we adopt the tool of slow SDE as in the aforementioned \citet{li2021happens} and \citet{gu2023and}.

% To fill in the gaps and provide a theoretical analysis that tracks iterations of Adam for a sufficiently long time, we use the same approximation tool as in the aforementioned \citet{li2021happens} and \citet{gu2023and}, namely \textit{slow SDE} (termed by \citet{gu2023and}). Specifically, optimizers such as SGD and Adam takes $\tilde{\O}(\eta^{-1})$ steps to converge onto the manifold, and then moves along the manifold for $\O(\eta^{-2})$ steps, during which the optimizer will be dominated by a slower implicit regularization dynamics, different from the mixing dynamics during the convergence phase. Conventional SDE papers such as \citet{li2018stochasticmodifiedequationsdynamics,li2021validitymodelingsgdstochastic,cattaneo2024implicitbiasadam,malladi2024sdesscalingrulesadaptive} approximate the iteration itself during the convergence phase, however, they struggle to bound the approximation error if extended to the manifold phase. 
% Instead, a slow SDE peels the convergence dynamics off by only approximating the iteration's projection on the manifold. In this way, a slow SDE can track the optimizer's iteration during the whole manifold phase for $\O(\eta^{-2})$ time. This idea will be made explicit in \Cref{sec:prelim}.

\paragraph{Adaptive Gradient Methods.}
As a test-of-time optimizer that has revolutionized the field of deep learning~\citep{kingma2014adam}, Adam innovatively combined the moving average of the first and second moments of gradients to determine an adaptive learning rate. Adam has also spawned a family of derivative optimizers such as AdamW \citep{loshchilov2017decoupled}, AdaFactor \citep{shazeer2018adafactor}, Adam-mini \citep{zhang2025adamminiusefewerlearning}, Adalayer \citep{zhao2025deconstructingmakesgoodoptimizer} and AdaSGD \citep{wang2020adasgd}, maintaining significant advantages over SGD in terms of empirical use. Under a more general framework of adaptive gradient methods, many optimizers also get huge success as adaptive gradient methods, such as RMSprop~\citep{hinton2012neural} and Adafactor~\citep{shazeer2018adafactor}.

\paragraph{Convergence of Adam.} There have been many previous works discussing the convergence bound of Adam. For example, \citet{reddi2018convergence} and \citet{dereich2024convergence} give convergence bounds under the convexity condition, \citet{zou2019sufficient}, \citet{shi2021rmsprop} and \citet{zhang2022adam} focus on the cases where learning rates follow a $1/\sqrt{t}$ decay, and the bounds given by \citet{NEURIPS2018_90365351}, \citet{zhang2022adam} and \citet{10.1145/3637528.3671718} do not decrease to $0$ as $\eta \to 0$. Also, most works \citep{defossez2020simple, guo2025unifiedconvergenceanalysisadaptive, iiduka2022theoreticalanalysisadamusing, wang2024convergenceadamnonuniformsmoothness, zhang2024convergence, hong2023highprobabilityconvergenceadam} only establish an upper bound on the average of gradient norms over the time of iteration. In this work, we derive a high-probability convergence analysis that directly bound the loss term of the last step, $\cL(\vtheta_k)-\cL^*$, to $o(1)$. Going beyond convex loss functions, we establish the bound on $\mu$-PL functions, and we focus on the constant learning rate schedule.